% =========================================================
% ARMIN OS -- TECHNICAL RESUME TEMPLATE
% Design: ThinkPad Cyber-Professional
% Format: A4 | Single Column | ATS Friendly
% =========================================================
<#
    This file is both the LaTeX design and a Jinja2 template that gets filled in
    from the site's own content (see portfolio_app/resume_latex.py). Editing the
    site is what changes the CV; there is no uploaded PDF to keep in step.

    The values use non-standard delimiters on purpose: LaTeX is full of braces
    and percent signs, so "\newcommand{\foo}[2]{%" would be read as a template
    tag and break the file. The delimiters are

        << value >>      a value from the database, already escaped for LaTeX
        <% if ... %>     a template tag: if / for / endif / endfor

    Everything outside those delimiters is ordinary LaTeX and passes through
    untouched, so the design below can be edited as normal.

    A value arrives already escaped, so never add an escaping or "safe" filter,
    and never remove the escaping in the module: a skill called "C++" or a
    description containing "50%" would stop the document compiling.

    The preamble works with both pdfLaTeX and XeTeX (Tectonic); see
    scripts/install-latex.sh for how the engine is obtained.
#>

\documentclass[a4paper,10pt]{article}

% ---------------------------------------------------------
% PACKAGES
% ---------------------------------------------------------

% iftex first: the input-encoding packages below are chosen by engine.
\usepackage{iftex}

\ifPDFTeX
  % pdfLaTeX reads bytes, so the encoding has to be declared, and
  % cmap/glyphtounicode are what make text copy-pasteable out of the PDF.
  % XeTeX and LuaTeX read UTF-8 natively and supply their own ToUnicode maps,
  % and would reject cmap.
  \usepackage[T1]{fontenc}
  \usepackage[utf8]{inputenc}
  \usepackage{cmap}
  \input{glyphtounicode}
  \pdfgentounicode=1
\fi

\usepackage{lmodern}
\usepackage{microtype}
\usepackage[empty]{fullpage}
\usepackage{titlesec}
\usepackage{xcolor}
\usepackage{enumitem}
\usepackage[hidelinks]{hyperref}
\usepackage{fancyhdr}
\usepackage[english]{babel}

% ---------------------------------------------------------
% ARMIN OS COLOR PALETTE
% ---------------------------------------------------------

\definecolor{cvgraphite}{HTML}{101820}
\definecolor{cvcyan}{HTML}{00A9D6}
\definecolor{cvdeepcyan}{HTML}{087EA4}
\definecolor{cvmuted}{HTML}{586776}
\definecolor{cvline}{HTML}{DCE4E9}

% ---------------------------------------------------------
% PDF METADATA
% ---------------------------------------------------------

\hypersetup{
    pdftitle={<< name >> - CV},
    pdfauthor={<< name >>},
    pdfsubject={<< title >>},
    pdfkeywords={<< keywords >>}
}

% ---------------------------------------------------------
% PAGE CONFIGURATION
% ---------------------------------------------------------

\pagestyle{fancy}
\fancyhf{}
\renewcommand{\headrulewidth}{0pt}
\renewcommand{\footrulewidth}{0pt}

\setlength{\footskip}{14pt}

\addtolength{\oddsidemargin}{-0.15in}
\addtolength{\evensidemargin}{-0.15in}
% The CV is one page, so the text block uses most of the sheet. Margins stay
% printable and the type size is untouched: it is the whitespace that gives, not
% the font.
\addtolength{\textwidth}{0.30in}

\addtolength{\topmargin}{-0.45in}
\addtolength{\textheight}{0.95in}

\urlstyle{same}

\raggedbottom
\raggedright
\setlength{\parindent}{0pt}
\setlength{\tabcolsep}{0pt}

% ---------------------------------------------------------
% SECTION STYLING
% ---------------------------------------------------------

\titleformat{\section}
  {\large\bfseries\color{cvgraphite}}
  {\color{cvcyan}\thesection}
  {0.65em}
  {}
  [\color{cvcyan}\titlerule]

% Spacing uses "x plus y minus z": the glue stretches and, crucially, shrinks.
% LaTeX therefore absorbs a small overflow by tightening the gaps instead of
% pushing the last few lines onto a second page, and spills over only when the
% content genuinely does not fit. The type size never changes.
\titlespacing*{\section}
  {0pt}
  {7pt plus 2pt minus 2pt}
  {4pt plus 1pt minus 1pt}

% ---------------------------------------------------------
% LIST CONFIGURATION
% ---------------------------------------------------------

\setlist[itemize]{
    leftmargin=1.4em,
    itemsep=1pt plus 1pt minus 0.5pt,
    topsep=2pt plus 1pt minus 1pt,
    parsep=0pt,
    partopsep=0pt
}

\renewcommand\labelitemii{\textbullet}

% ---------------------------------------------------------
% CUSTOM COMMANDS
% ---------------------------------------------------------

% Hyperlinks
\newcommand{\resumelink}[2]{%
    \href{#1}{\textcolor{cvdeepcyan}{#2}}%
}

% Experience / Education heading
\newcommand{\resumeSubheading}[4]{%
    \item
    \textbf{#1}\hfill
    {\small\color{cvmuted}#2}\\
    {\small #3}\hfill
    {\small\color{cvmuted}#4}
    \vspace{-2pt}
}

% Project heading
\newcommand{\resumeProjectHeading}[2]{%
    \item
    \textbf{#1}\hfill
    {\small\color{cvmuted}#2}
    \vspace{-2pt}
}

% Description bullet
\newcommand{\resumeItem}[1]{%
    \item {\small #1}
}

% List wrappers
\newcommand{\resumeSubHeadingListStart}{
    \begin{itemize}[leftmargin=0pt,label={},itemsep=3pt plus 1pt minus 1.5pt]
}

\newcommand{\resumeSubHeadingListEnd}{
    \end{itemize}
}

\newcommand{\resumeItemListStart}{
    \begin{itemize}
}

\newcommand{\resumeItemListEnd}{
    \end{itemize}\vspace{-3pt}
}

% ---------------------------------------------------------
% DOCUMENT START
% ---------------------------------------------------------

\begin{document}

% =========================================================
% HEADER
% =========================================================

\begin{center}

    {\fontsize{21}{24}\selectfont
    \bfseries\color{cvgraphite}
    << name >>}

    \vspace{3pt}

    {\small\bfseries\color{cvdeepcyan}
    << title >>}

    \vspace{5pt}

    {\small
    <% for item in contact %><% if item.url %>\resumelink{<< item.url >>}{<< item.label >>}<% else %><< item.label >><% endif %><% if not loop.last %> \;|\; <% endif %><% endfor %>}

\end{center}

\vspace{-4pt}
{\color{cvcyan}\hrule height 1.2pt}
\vspace{3pt}

% =========================================================
% PROFESSIONAL SUMMARY
% =========================================================

<% if summary %>
\section{Professional Summary}

\small
<< summary >>
\normalsize
<% endif %>

% =========================================================
% EXPERIENCE
% =========================================================

<% if experiences %>
\section{Experience}

\resumeSubHeadingListStart
<% for experience in experiences %>
    \resumeSubheading
        {<< experience.company >>}
        {<< experience.dates >>}
        {<< experience.role >>}
        {<< experience.location >>}
<% if experience.bullets %>
    \resumeItemListStart
<% for bullet in experience.bullets %>
        \resumeItem{<< bullet >>}
<% endfor %>
    \resumeItemListEnd
<% endif %>
<% endfor %>
\resumeSubHeadingListEnd
<% endif %>

% =========================================================
% EDUCATION
% =========================================================

<% if educations %>
\section{Education}

\resumeSubHeadingListStart
<% for education in educations %>
    \resumeSubheading
        {<< education.institution >>}
        {<< education.dates >>}
        {<< education.degree >>}
        {<< education.location >>}
<% if education.bullets %>
    \resumeItemListStart
<% for bullet in education.bullets %>
        \resumeItem{<< bullet >>}
<% endfor %>
    \resumeItemListEnd
<% endif %>
<% endfor %>
\resumeSubHeadingListEnd
<% endif %>

% =========================================================
% SELECTED PROJECTS
% =========================================================

<% if projects %>
\section{Selected Projects}

\resumeSubHeadingListStart
<% for project in projects %>
    \resumeProjectHeading
        {\resumelink{<< project.link >>}{<< project.name >><% if project.technologies %> $|$ << project.technologies >><% endif %>}}
        {}
<% if project.bullets %>
    \resumeItemListStart
<% for bullet in project.bullets %>
        \resumeItem{<< bullet >>}
<% endfor %>
    \resumeItemListEnd
<% endif %>
<% endfor %>
\resumeSubHeadingListEnd
<% endif %>

% =========================================================
% TECHNICAL SKILLS
% =========================================================

% One line per main skill, with that skill's sub-skills as the detail. The
% two-level skill tree on the site already groups technologies this way, so no
% separate categorisation is needed here.

<% if skills %>
\section{Technical Skills}

\begin{itemize}[leftmargin=0pt,label={},itemsep=1.5pt plus 1pt minus 0.5pt]

<% for skill in skills %>
    \item \small{
        \textbf{<< skill.name >><% if skill.technologies %>: << skill.technologies >><% endif %>}
    }
<% endfor %>

\end{itemize}
<% endif %>

% =========================================================
% ADDITIONAL INFORMATION
% =========================================================

<% if languages or certificates %>
\section{Additional Information}

\begin{itemize}[leftmargin=0pt,label={},itemsep=1.5pt plus 1pt minus 0.5pt]

<% if languages %>
    \item \small{
        \textbf{Languages:}
        <% for language in languages %><< language.name >> (<< language.level >>)<% if not loop.last %>, <% endif %><% endfor %>
    }
<% endif %>

<% if certificates %>
    \item \small{
        \textbf{Certifications:}
        <% for certificate in certificates %><< certificate.name >> (<< certificate.institution >>, << certificate.year >>)<% if not loop.last %>; <% endif %><% endfor %>
    }
<% endif %>

\end{itemize}
<% endif %>

% =========================================================
% DOCUMENT END
% =========================================================

\end{document}
